\chapter{Giới thiệu chung}

% ==========================================================
\section{Tổng quan về hệ thống AVI trong công nghiệp}

Trong bối cảnh cạnh tranh sản xuất ngày càng khốc liệt, việc duy trì chất lượng sản phẩm ổn định không chỉ là lợi thế mà đã trở thành yêu cầu bắt buộc. Các lỗi ngoại quan, dù rất nhỏ, đều có thể gây ra chi phí lớn như thu hồi sản phẩm, gián đoạn dây chuyền và ảnh hưởng nghiêm trọng đến uy tín thương hiệu. Kiểm tra ngoại quan thủ công, vốn phụ thuộc nhiều vào con người, khó đảm bảo tính ổn định và tốc độ trong dây chuyền hiện đại.

Hệ thống kiểm tra ngoại quan tự động — \textit{Automated Visual Inspection} (AVI) — ra đời như một giải pháp tất yếu, tận dụng thị giác máy và trí tuệ nhân tạo để tự động phát hiện lỗi và đánh giá chất lượng theo thời gian thực. Theo dự báo, thị trường thị giác máy toàn cầu sẽ đạt khoảng \textbf{9,29 tỷ USD} vào năm 2032 với tốc độ tăng trưởng 7,2\%/năm \cite{easyodm_AVI2025}, phản ánh nhu cầu rất lớn về tự động hoá và kiểm tra chất lượng trong Công nghiệp 4.0.

\subsection*{Khái niệm và chức năng của hệ thống AVI}

AVI là hệ thống kiểm tra chất lượng dựa trên camera công nghiệp, hệ thống chiếu sáng và thuật toán xử lý ảnh. Các chức năng chính bao gồm:

\begin{itemize}
    \item Phát hiện và phân loại lỗi trên sản phẩm.
    \item Định vị lỗi và đánh giá độ nghiêm trọng.
    \item Nhận diện bất thường (\textit{anomaly detection}) trên bề mặt.
    \item Đánh giá chất lượng tổng thể và đưa ra quyết định theo thời gian thực.
\end{itemize}

Các thuật toán phổ biến trong AVI gồm có Machine Learning, Deep Learning (YOLO, CNN), phương pháp phát hiện bất thường, trích xuất đặc trưng và mô hình so khớp hình ảnh.

\subsection*{Quy trình hoạt động của hệ thống AVI}

Quy trình điển hình của một hệ thống AVI bao gồm:

\begin{enumerate}
    \item \textbf{Thu nhận hình ảnh}: camera độ phân giải cao kết hợp ánh sáng phù hợp tạo ra ảnh sắc nét.
    \item \textbf{Xử lý ảnh}: lọc nhiễu, cân bằng sáng, phân đoạn và trích xuất đặc trưng.
    \item \textbf{So sánh và phân tích}: đối chiếu với tiêu chuẩn hoặc mô hình AI đã học.
    \item \textbf{Ra quyết định}: phân loại sản phẩm đạt (OK) hoặc lỗi (NG).
    \item \textbf{Phản hồi}: truyền tín hiệu điều khiển để loại bỏ sản phẩm lỗi hoặc ghi log dữ liệu.
\end{enumerate}

\subsection*{Các loại lỗi hệ thống AVI có thể phát hiện}

AVI có khả năng phát hiện nhiều loại lỗi ngoại quan:

\begin{itemize}
    \item \textbf{Lỗi bề mặt}: trầy xước, móp, đổi màu, nứt, ba via.
    \item \textbf{Lỗi kích thước}: sai lệch hình dạng, thể tích, vị trí.
    \item \textbf{Lỗi lắp ráp}: thiếu linh kiện, sai vị trí, lỗi hàn.
    \item \textbf{Lỗi bao bì}: nhãn sai, seal lỗi, mức đầy không đúng.
    \item \textbf{Lỗi vật liệu}: tạp chất, lẫn bụi, rỗ khí.
\end{itemize}


% ==========================================================
\section{Kiến trúc và thành phần của hệ thống AVI}

Một hệ thống AVI hoàn chỉnh được xây dựng từ sự kết hợp của phần cứng chuyên dụng và phần mềm xử lý mạnh mẽ nhằm đảm bảo tốc độ và độ chính xác kiểm tra.

\subsection*{Thành phần phần cứng}

\begin{itemize}
    \item \textbf{Camera công nghiệp}: area-scan, line-scan, camera 3D, multispectral, hoặc X-ray tùy ứng dụng.
    \item \textbf{Hệ thống chiếu sáng}: ring light, dome light, backlight, directional lighting, nhằm làm nổi bật đặc trưng bề mặt.
    \item \textbf{Ống kính và quang học}: lens tiêu chuẩn hoặc telecentric cho bài toán đo kích thước chính xác cao.
    \item \textbf{Cảm biến hỗ trợ}: cảm biến 3D, cảm biến tiệm cận, IR hoặc LIDAR giúp bổ sung thông tin ngoài hình ảnh.
    \item \textbf{Máy tính xử lý}: CPU/GPU hiệu năng cao để chạy các mô hình AI và xử lý ảnh thời gian thực.
\end{itemize}

\subsection*{Thành phần phần mềm}

\begin{itemize}
    \item \textbf{Xử lý ảnh}: bao gồm tiền xử lý, lọc nhiễu, phân đoạn, trích xuất đặc trưng.
    \item \textbf{AI và học máy}: YOLO, CNN, anomaly detection dùng để phát hiện lỗi phức tạp.
    \item \textbf{Giao diện người dùng (UI)}: hiển thị ảnh, kết quả kiểm tra và thống kê.
    \item \textbf{Kết nối mạng}: Ethernet/IP, OPC UA, MQTT để tích hợp vào dây chuyền.
    \item \textbf{Lưu trữ dữ liệu}: lưu ảnh kiểm tra, log sự kiện và thông số phục vụ phân tích chất lượng.
\end{itemize}


% ==========================================================
\section{Ứng dụng, lợi ích, thách thức và xu hướng phát triển}

\subsection*{Ứng dụng của AVI trong công nghiệp}

\begin{itemize}
    \item \textbf{Ô tô}: kiểm tra sơn, mối hàn, lắp ráp linh kiện.
    \item \textbf{Điện tử}: phát hiện lỗi PCB, kiểm tra mạch SMT, linh kiện thiếu hoặc sai loại.
    \item \textbf{Dược phẩm}: kiểm tra bao bì, nhãn, dị vật.
    \item \textbf{Thực phẩm}: kiểm tra mức đầy, lỗi bao bì, dị vật.
    \item \textbf{Gỗ, dệt may, bán dẫn, hàng không}: phát hiện lỗi bề mặt, nứt, sai kích thước.
\end{itemize}

\subsection*{Lợi ích nổi bật}

\begin{itemize}
    \item Độ chính xác cao (95--99,5\%), ổn định hơn kiểm tra thủ công.
    \item Tốc độ xử lý nhanh (0,1–0,5 giây/sản phẩm), vận hành 24/7.
    \item Giảm chi phí nhân lực, giảm phế phẩm, thời gian hoàn vốn ngắn (12--24 tháng).
    \item Nâng cao chất lượng đầu ra và tính nhất quán.
    \item Tích hợp linh hoạt vào dây chuyền, dễ mở rộng.
\end{itemize}

\subsection*{Thách thức và hạn chế}

\begin{itemize}
    \item Chi phí đầu tư ban đầu tương đối cao.
    \item Nhạy cảm với ánh sáng và vị trí sản phẩm.
    \item Khó phát hiện lỗi cực nhỏ nếu không có giải pháp ánh sáng phù hợp.
    \item Hệ thống AI yêu cầu dữ liệu huấn luyện đa dạng.
    \item Tích hợp vào dây chuyền hiện tại đôi khi phức tạp.
\end{itemize}

\subsection*{Xu hướng phát triển}

\begin{itemize}
    \item Tích hợp AI nâng cao và mô hình tự học.
    \item Hình ảnh 3D, hyperspectral, X-ray cho khả năng phân tích sâu hơn.
    \item Edge computing giúp xử lý thời gian thực với độ trễ thấp.
    \item Robot cộng tác và thị giác máy kết hợp cho kiểm tra linh hoạt đa góc.
    \item Hệ thống no-code cho phép kỹ thuật viên cấu hình nhanh mà không cần lập trình.
\end{itemize}

\subsection*{Kết luận}

AVI đang trở thành một thành phần không thể thiếu trong các dây chuyền sản xuất hiện đại, giúp doanh nghiệp nâng cao chất lượng sản phẩm, tối ưu chi phí và duy trì lợi thế cạnh tranh. Sự phát triển mạnh mẽ của AI và thị giác máy dự báo rằng vai trò của AVI sẽ ngày càng quan trọng trong tương lai của sản xuất thông minh.
